%%% FIELD proof-universal-validity-iff
Fix \(s\in\mathbb R^{\Omega}\) and \(\alpha\in[0,1]\), and define
\[
R=\{x\in\Omega:p_K(x;s)\leq\alpha\}.
\]
If \(R=\varnothing\), then
\(\Pr_{X\sim\pi}\{p_K(X;s)\leq\alpha\}=0\leq\alpha\).  Suppose
\(R\neq\varnothing\).  Because \(R\) is finite, it has an ordering
\(\sigma_1,\ldots,\sigma_{|R|}\) satisfying
\[
s(\sigma_1)\geq\cdots\geq s(\sigma_{|R|});
\]
ties may be broken arbitrarily.  For
\(P_j=\{\sigma_1,\ldots,\sigma_j\}\), this ordering gives
\[
P_j\subseteq\{y\in\Omega:s(y)\geq s(\sigma_j)\}.
\tag{1}
\]
The candidate-kernel type gives \(K(x,y)\geq0\), while
\(\texttt{ass:row-stochastic}\) gives unit row sums.  Hence the row masses are
monotone in their set argument.  By (1) and
\(\texttt{def:kernel-tail-value}\), for every \(j\),
\[
K(\sigma_j,P_j)
\leq K\!\left(\sigma_j,\{y:s(y)\geq s(\sigma_j)\}\right)
=p_K(\sigma_j;s)
\leq\alpha.
\tag{2}
\]
Taking the maximum in (2) and then applying
\(\texttt{def:subset-degeneracy}\) yields
\[
d_K(R)
\leq\max_{1\leq j\leq |R|}K(\sigma_j,P_j)
\leq\alpha.
\tag{3}
\]
If \(d_K(A)\geq\pi(A)\) for every nonempty \(A\), then (3) and the design
law \(X\sim\pi\) give
\[
\Pr_{X\sim\pi}\{p_K(X;s)\leq\alpha\}
=\pi(R)
\leq d_K(R)
\leq\alpha.
\]
Since \(s\) and \(\alpha\) were arbitrary, this is universal validity by
\(\texttt{def:universal-validity}\).

Conversely, suppose that the subset inequality fails at a nonempty
\(A\subseteq\Omega\), so
\[
d_K(A)<\pi(A).
\tag{4}
\]
The set \(\operatorname{Ord}(A)\) is finite, so the minimum in
\(\texttt{def:subset-degeneracy}\) is attained by some ordering
\(\sigma_1,\ldots,\sigma_{|A|}\).  Choose a score with
\[
s(\sigma_1)>\cdots>s(\sigma_{|A|})
\]
and assign every point in \(\Omega\setminus A\) a score strictly below
\(s(\sigma_{|A|})\).  For this score, the upper score set at \(\sigma_j\) is
exactly \(P_j\).  Therefore
\[
p_K(\sigma_j;s)
=K(\sigma_j,P_j)
\leq\max_{1\leq i\leq |A|}K(\sigma_i,P_i)
=d_K(A),
\qquad j=1,\ldots,|A|.
\tag{5}
\]
Row stochasticity ensures \(d_K(A)\in[0,1]\), so
\(\alpha=d_K(A)\) is an admissible threshold.  Equation (5) shows that the
rejection set at this threshold contains \(A\).  Using (4),
\[
\Pr_{X\sim\pi}\{p_K(X;s)\leq d_K(A)\}
\geq\pi(A)
>d_K(A),
\]
which violates \(\texttt{def:universal-validity}\).  Thus universal validity
implies every subset inequality.  The use of \(\leq\) in (5) also proves that
there is no boundary loss when a tail value equals the rejection threshold.

To make the tightening relative to partition sufficiency explicit, let
\(\mathcal G\) be any partition of \(\Omega\), write \(G_x\) for the cell
containing \(x\), and consider the conditional row kernel
\[
K(x,B)=\frac{\pi(B\cap G_x)}{\pi(G_x)}.
\tag{6}
\]
The denominator is positive by \(\texttt{ass:positive-finite-design}\).  Fix a
nonempty \(A\) and any ordering of it.  For each cell \(G\) meeting \(A\), look
at the last element of \(A\cap G\) in that ordering.  At its prefix, the row
mass is \(\pi(A\cap G)/\pi(G)\).  Consequently the maximum prefix mass is at
least
\[
\max_{G:\,A\cap G\neq\varnothing}
\frac{\pi(A\cap G)}{\pi(G)}
\geq
\sum_{G\in\mathcal G}\pi(G)
\frac{\pi(A\cap G)}{\pi(G)}
=\pi(A),
\tag{7}
\]
where cells disjoint from \(A\) contribute zero and the middle inequality says
that a maximum dominates a weighted average.  Since the ordering was arbitrary,
\(d_K(A)\geq\pi(A)\).  Thus every partition kernel is included in the
characterized class.

It remains to verify that the inclusion is strict.  Let
\(\Omega=\{1,2,3\}\), let \(\pi\) be uniform, and, with indices interpreted
cyclically, set
\[
K(i,i)=\frac23,
\qquad
K(i,i+1)=\frac13,
\qquad
K(i,j)=0\ \text{otherwise}.
\tag{8}
\]
The matrix in (8) is nonnegative and row-stochastic.  For a singleton,
\(d_K(A)=2/3\).  Any two-point set contains exactly one directed cycle edge.
Order its tail first and its head second.  The first prefix contains only the
tail, while the head's outgoing cycle edge leaves the set, so both prefix
masses are \(2/3\).  Every ordering has first-prefix mass \(2/3\).  Hence
\(d_K(A)=2/3\) for every two-point set.  Finally,
\(d_K(\Omega)=1\), because the last prefix in every ordering is \(\Omega\) and
has row mass one.  Thus
\[
d_K(A)\geq\pi(A)
\qquad(\varnothing\neq A\subseteq\Omega),
\]
and the equivalence just proved makes (8) universally valid.  Taking the legal
relation to contain the diagonal and the three directed cycle edges makes (8)
a legal kernel.  Its three two-point row supports overlap but are neither
identical nor disjoint.  By contrast, rows obtained by conditioning \(\pi\) on
partition cells have identical supports within a cell and disjoint supports
across cells.  Therefore (8) is not a partition kernel, proving that the
characterized class admits legal nonpartition kernels.

%%% FIELD construction-worst-case-change
For every observed \(x_0\in\Omega\),
\[
p^{\mathrm{wc}}_{r,c}(x_0)
=\sup_{\theta\in\Theta_{r,c}(x_0,Y^{\mathrm{obs}})}
p_K(x_0;s_{\theta}).
\]
For this worst-case p-value, set \(\sup\varnothing=0\).  Thus
\(p^{\mathrm{wc}}_{r,c}(x_0)\in[0,1]\) even when no compatible schedule
exists.

%%% FIELD construction-worst-case-core
For every observed \(x_0\in\Omega\),
\[
p^{\mathrm{wc}}_{r,c}(x_0)
=\sup_{\theta\in\Theta_{r,c}(x_0,Y^{\mathrm{obs}})}
p_K(x_0;s_{\theta}).
\]
For this worst-case p-value, set \(\sup\varnothing=0\).  Thus
\(p^{\mathrm{wc}}_{r,c}(x_0)\in[0,1]\) even when no compatible schedule
exists.

%%% FIELD construction-joint-region-change
The joint inverted schedule region is
\[
\mathcal C_{\alpha}
=\bigcap_{r=0}^{N}
  \bigcap_{c:\,p^{\mathrm{wc}}_{r,c}(X)\leq\alpha}
  \left\{\theta:
  \sum_{u\in U}\mathbf 1\{\theta_u(1,e)-\theta_u(0,e)>c\}>r
  \right\}.
\]
An inner intersection indexed by no threshold \(c\) is the full schedule space
\(\mathbb R^{U\times\{0,1\}\times\mathcal E}\).  A rejected constraint with
\(r=N\) contributes the empty set because a sum of \(N\) indicators cannot
exceed \(N\).

%%% FIELD construction-joint-region-core
The joint inverted schedule region is
\[
\mathcal C_{\alpha}
=\bigcap_{r=0}^{N}
  \bigcap_{c:\,p^{\mathrm{wc}}_{r,c}(X)\leq\alpha}
  \left\{\theta:
  \sum_{u\in U}\mathbf 1\{\theta_u(1,e)-\theta_u(0,e)>c\}>r
  \right\}.
\]
An inner intersection indexed by no threshold \(c\) is the full schedule space
\(\mathbb R^{U\times\{0,1\}\times\mathcal E}\).  A rejected constraint with
\(r=N\) contributes the empty set because a sum of \(N\) indicators cannot
exceed \(N\).

%%% FIELD construction-quantile-bound-change
For \(1\leq k\leq N\), define
\[
\ell_k(e;\alpha)
=\inf_{\theta\in\mathcal C_{\alpha}}\theta_{(k)}(e),
\qquad
\inf\varnothing=+\infty.
\tag{1}
\]
The following clause evaluates (1).  For an observed \(x_0\in\Omega\), set
\(r=N-k\).  For each \(c\in\mathbb R\), enumerate binary labels
\(h\in\{0,1\}^{U}\) with \(\sum_{u\in U}h_u\leq r\) and ternary
assignment-wide weak-order labels
\[
\lambda_{xuv}\in\{-1,0,1\},
\qquad x\in\Omega,\quad 1\leq u<v\leq N.
\]
Writing \(z_{ua}=\theta_u(a,e)\), impose
\[
z_{u,Z_u(x_0)}=Y_u^{\mathrm{obs}},
\qquad
h_u=1\Longrightarrow z_{u1}-z_{u0}>c,
\qquad
h_u=0\Longrightarrow z_{u1}-z_{u0}\leq c,
\]
and let \(\lambda_{xuv}\) record the sign of
\(z_{u,Z_u(x)}-z_{v,Z_v(x)}\).  These are the observation, effect, and
weak-order constraints for the label.
After writing its strict rows as \(a_j^{\mathsf T}z<b_j\), its weak rows as
\(Cz\leq d\), and its equalities as \(Ez=f\), test a cell containing strict
rows with
\[
\max_{z,\varepsilon}\ \varepsilon
\quad\text{subject to}\quad
a_j^{\mathsf T}z\leq b_j-\varepsilon\ (j),
\quad Cz\leq d,
\quad Ez=f,
\quad 0\leq\varepsilon\leq1;
\tag{2}
\]
the cell is feasible exactly when the optimum in (2) is positive.  If there is
no strict row, use ordinary linear feasibility.  The fixed tie rule assigns to
each feasible \(\lambda\) the constants
\(t_{\lambda}(x)=T(x,Y_{\theta}(x))\), and
\[
p^{\mathrm{wc}}_{r,c}(x_0)
=\max_{(h,\lambda):\,\mathrm{feasible}}
\sum_{y\in\Omega}K(x_0,y)
\mathbf 1\{t_{\lambda}(y)\geq t_{\lambda}(x_0)\},
\qquad \max\varnothing=0.
\tag{3}
\]
Define
\[
b_k(x_0;\alpha)
=\sup\left\{c\in\mathbb R:
p^{\mathrm{wc}}_{N-k,c}(x_0)\leq\alpha\right\},
\qquad \sup\varnothing=-\infty.
\tag{4}
\]
The projected bound remains defined by (1), not by the computation below.  The
exact network schedule theorem proves the evaluation identity
\[
\ell_k(e;\alpha)=b_k(X;\alpha).
\tag{5}
\]
Treating \(c\) as a variable, projection of every feasible weak-order cell onto
the \(c\)-line gives an interval with its endpoint inclusions recorded by (2).
For rational or real-algebraic representations of the kernel weights, observed
outcomes, level, and finite rank-score table, the finitely many projected
intervals and endpoints give a terminating exact evaluation of (3)--(5).  With
arbitrary unspecified real inputs, (2)--(5) retain only their finite ideal
exact-real interpretation.

%%% FIELD construction-quantile-bound-core
For \(1\leq k\leq N\), define
\[
\ell_k(e;\alpha)
=\inf_{\theta\in\mathcal C_{\alpha}}\theta_{(k)}(e),
\qquad
\inf\varnothing=+\infty.
\tag{1}
\]
The following clause evaluates (1).  For an observed \(x_0\in\Omega\), set
\(r=N-k\).  For each \(c\in\mathbb R\), enumerate binary labels
\(h\in\{0,1\}^{U}\) with \(\sum_{u\in U}h_u\leq r\) and ternary
assignment-wide weak-order labels
\[
\lambda_{xuv}\in\{-1,0,1\},
\qquad x\in\Omega,\quad 1\leq u<v\leq N.
\]
Writing \(z_{ua}=\theta_u(a,e)\), impose
\[
z_{u,Z_u(x_0)}=Y_u^{\mathrm{obs}},
\qquad
h_u=1\Longrightarrow z_{u1}-z_{u0}>c,
\qquad
h_u=0\Longrightarrow z_{u1}-z_{u0}\leq c,
\]
and let \(\lambda_{xuv}\) record the sign of
\(z_{u,Z_u(x)}-z_{v,Z_v(x)}\).  These are the observation, effect, and
weak-order constraints for the label.
After writing its strict rows as \(a_j^{\mathsf T}z<b_j\), its weak rows as
\(Cz\leq d\), and its equalities as \(Ez=f\), test a cell containing strict
rows with
\[
\max_{z,\varepsilon}\ \varepsilon
\quad\text{subject to}\quad
a_j^{\mathsf T}z\leq b_j-\varepsilon\ (j),
\quad Cz\leq d,
\quad Ez=f,
\quad 0\leq\varepsilon\leq1;
\tag{2}
\]
the cell is feasible exactly when the optimum in (2) is positive.  If there is
no strict row, use ordinary linear feasibility.  The fixed tie rule assigns to
each feasible \(\lambda\) the constants
\(t_{\lambda}(x)=T(x,Y_{\theta}(x))\), and
\[
p^{\mathrm{wc}}_{r,c}(x_0)
=\max_{(h,\lambda):\,\mathrm{feasible}}
\sum_{y\in\Omega}K(x_0,y)
\mathbf 1\{t_{\lambda}(y)\geq t_{\lambda}(x_0)\},
\qquad \max\varnothing=0.
\tag{3}
\]
Define
\[
b_k(x_0;\alpha)
=\sup\left\{c\in\mathbb R:
p^{\mathrm{wc}}_{N-k,c}(x_0)\leq\alpha\right\},
\qquad \sup\varnothing=-\infty.
\tag{4}
\]
The projected bound remains defined by (1), not by the computation below.  The
exact network schedule theorem proves the evaluation identity
\[
\ell_k(e;\alpha)=b_k(X;\alpha).
\tag{5}
\]
Treating \(c\) as a variable, projection of every feasible weak-order cell onto
the \(c\)-line gives an interval with its endpoint inclusions recorded by (2).
For rational or real-algebraic representations of the kernel weights, observed
outcomes, level, and finite rank-score table, the finitely many projected
intervals and endpoints give a terminating exact evaluation of (3)--(5).  With
arbitrary unspecified real inputs, (2)--(5) retain only their finite ideal
exact-real interpretation.
